Une ampoule incandescente de \qty{\v{P1}}{W} a un rendement lumineux de \qty{\v{e1}}{\percent}. On la remplace par une ampoule DEL de \qty{\v{P2}}{W} qui produit exactement la même quantité de lumière. L'ampoule est allumée \qty{\v{h}}{h} par jour, \qty{365}{jours} par année.

\begin{enumerate}[label=\alph*)]
    \item Quelle puissance lumineuse (énergie lumineuse par seconde) l'ampoule incandescente produit-elle ? En quoi est transformé le reste de l'énergie ?
    \item Calculez le rendement de l'ampoule DEL.
    \item Quelle quantité d'énergie, en kilowattheures, économise-t-on en une année avec la DEL ?
\end{enumerate}
